\section*{Chapitre \ref{ch:b2:realfun} --- Fonctions d'une variable réelle}

\begin{solution}{exo:b2:realfun:1}
$\lfloor x\rfloor$~: sauts de taille $1$ en tout entier. $x -
\lfloor x\rfloor$~: sauts de taille $-1$ aux entiers (limite à
gauche $1$, valeur $0$). $\lfloor x\rfloor + \sqrt{x - \lfloor
x\rfloor}$~: en un entier $n$, limite à gauche $(n - 1) + 1 = n$ et
valeur $n$~: \emph{\omterm{def:b2:metric:continuity}{continue}} partout (la racine carrée répare le
saut), bien que non dérivable aux entiers. La fonction à sauts
rationnels~: en restreignant la construction à une énumération des
dyadiques, elle saute de $2^{-n}$ exactement au $n$-ième rationnel
dyadique et est \omterm{def:b2:metric:continuity}{continue} ailleurs.
\end{solution}

\begin{solution}{exo:b2:realfun:2}
Supposons que $f$ croissante ait une discontinuité en un point
intérieur $a$~: alors $f(a^-) < f(a^+)$
(le~\cref{thm:b2:realfun:monotone}) et l'image de $I$ évite
l'intervalle \omterm{def:b2:metric:topology}{ouvert} non vide $\intoo{f(a^-)}{f(a^+)}$ sauf peut-être
la seule valeur $f(a)$~: l'image de tout sous-intervalle contenant
$a$ dans son intérieur n'est pas un intervalle (il a un trou d'au
moins un côté de $f(a)$). Cela contredit la propriété des valeurs
intermédiaires. Les discontinuités aux extrémités sont exclues de la
même manière avec des trous unilatéraux.
\end{solution}

\begin{solution}{exo:b2:realfun:3}
$x\ln x$~: dérivée seconde $\frac1x > 0$~: convexe. $\ln(1 +
\eu^x)$~: dérivée $\frac{\eu^x}{1 + \eu^x} = 1 - \frac{1}{1 +
\eu^x}$, croissante~: convexe. $\sqrt{1 + x^2}$~: dérivée seconde
$(1 + x^2)^{-3/2} > 0$~: convexe. $x^3$~: non convexe sur $\R$ ($f''
= 6x$ change de signe)~; convexe seulement sur $\R_+$.
\end{solution}

\begin{solution}{exo:b2:realfun:4}
Supposons $f(a) \neq f(b)$, disons $f(b) > f(a)$ avec $a < b$ (le cas
$f(b) < f(a)$ est symétrique, en regardant à gauche). Pour $x > b$,
l'inégalité des pentes (le~\cref{lem:b2:realfun:slopes}) sur $a < b <
x$ donne
\[
\frac{f(x) - f(a)}{x - a} \geq \frac{f(b) - f(a)}{b - a} = m > 0
\quad\Longrightarrow\quad
f(x) \geq f(a) + m(x - a) \xrightarrow[x\to+\infty]{} +\infty,
\]
ce qui contredit le caractère majoré. Donc $f$ est constante.

Asymptotes~: si $f(x) - (\alpha x + \beta) \to 0$ en $+\infty$ et
$f(x) - (\alpha' x + \beta') \to 0$ en $-\infty$, la fonction convexe
$g(x) = f(x) - (\alpha x + \beta)$ est majorée au voisinage de
$+\infty$~; la convexité plus une asymptote en $-\infty$ (qui force
$\alpha' \leq \alpha$ puis $\alpha' = \alpha$ en comparant les pentes
en $\mp\infty$~: les pentes d'une fonction convexe croissent) rend
$g$ majorée sur tout $\R$, donc constante $= 0$ à la limite~: $f$ est
affine.
\end{solution}

\begin{solution}{exo:b2:realfun:5}
$f'$ est monotone, donc par le~\cref{thm:b2:realfun:monotone} ses
seules discontinuités possibles sont des sauts, avec des limites
latérales existant partout. Par le~\cref{cor:b2:realfun:nojumps}, une
dérivée n'a pas de discontinuité de saut. Donc $f'$ n'a aucune
discontinuité~: \omterm{def:b2:metric:continuity}{continue}.
\end{solution}

\begin{solution}{exo:b2:realfun:6}
Prenons les logarithmes~:
\[
\frac1p \ln\Bigl(\sum_i \lambda_i x_i^p\Bigr)
= \frac1p \ln\Bigl(1 + p\sum_i \lambda_i \ln x_i +
O(p^2)\Bigr)
= \sum_i \lambda_i \ln x_i + O(p)
\xrightarrow[p \to 0^+]{} \sum_i \lambda_i \ln x_i ,
\]
en utilisant $\sum\lambda_i = 1$ et $\ln(1 + u) = u + O(u^2)$.
En exponentiant on obtient la moyenne géométrique. Or pour tout $p
\in \intoo{0}{1}$, l'inégalité des moyennes de puissances
(l'\cref{ex:b2:realfun:powermeans}, exposants $p < 1$) donne
\[
\Bigl(\sum_i \lambda_i x_i^p\Bigr)^{1/p} \leq \sum_i\lambda_i x_i ;
\]
en faisant $p \to 0^+$ à gauche on obtient $\prod_i x_i^{\lambda_i}
\leq \sum_i \lambda_i x_i$~: l'inégalité arithmético-géométrique
pondérée.
\end{solution}

\begin{solution}{exo:b2:realfun:7}
Avec $\varphi(t) = t\ln t$ (convexe~: $\varphi'' = \frac1t > 0$) et
les poids $q_i$ aux points $t_i = \frac{p_i}{q_i}$~:
\[
\sum_i p_i \ln\frac{p_i}{q_i}
= \sum_i q_i\, \varphi\Bigl(\frac{p_i}{q_i}\Bigr)
\;\geq\; \varphi\Bigl(\sum_i q_i \frac{p_i}{q_i}\Bigr)
= \varphi(1) = 0 ,
\]
par Jensen (le~\cref{thm:b2:realfun:convexreg}~(3)). L'égalité dans
Jensen pour une fonction \emph{strictement} convexe force tous les
points $t_i$ à coïncider~: $\frac{p_i}{q_i}$ constant, et en sommant,
la constante est $1$~: $p = q$. (Cette quantité --- la divergence de
Kullback--Leibler --- revient dans le monde du~\cref{ch:b2:randomvar}.)
\end{solution}

\begin{solution}{exo:b2:realfun:8}
\emph{Poids dyadiques.} Par récurrence sur $m$~: le cas $m = 1$ est
l'hypothèse. Pour un poids $\lambda = \frac{k}{2^{m+1}}$ ($k$
impair), écrivons $\lambda = \frac12(\lambda_1 + \lambda_2)$ avec
$\lambda_j = \frac{k \mp 1}{2^{m+1}}$, tous deux de dénominateur
$2^m$ après simplification~; alors
\[
f\bigl(\lambda x + (1{-}\lambda)y\bigr)
= f\Bigl(\tfrac{u + v}{2}\Bigr)
\leq \frac{f(u) + f(v)}{2}
\leq \lambda f(x) + (1 - \lambda) f(y),
\]
où $u = \lambda_1 x + (1 - \lambda_1)y$ et $v = \lambda_2 x +
(1-\lambda_2)y$, en utilisant la convexité au milieu puis
l'hypothèse de récurrence sur $u, v$.

\emph{Passage à la limite.} Pour un $\lambda \in \intcc{0}{1}$
quelconque, prenons des dyadiques $\lambda_n \to \lambda$~: la
continuité de $f$ et des applications affines passe l'inégalité
$f(\lambda_n x + (1-\lambda_n)y) \leq \lambda_n f(x) +
(1-\lambda_n)f(y)$ à la limite~: $f$ est convexe.
\end{solution}

\begin{solution}{exo:b2:realfun:9}
Pour $0 < x < y$~: l'inégalité des pentes
(le~\cref{lem:b2:realfun:slopes}) aux points $0 < x < y$ donne
\[
\frac{f(x) - f(0)}{x} \leq \frac{f(y) - f(0)}{y},
\qquad\text{c.-à-d.}\qquad
\frac{f(x)}{x} \leq \frac{f(y)}{y} +
f(0)\Bigl(\frac1x - \frac1y\Bigr).
\]
Comme $f(0) \leq 0$ et $\frac1x - \frac1y > 0$, le dernier terme est
$\leq 0$~: $\frac{f(x)}{x} \leq \frac{f(y)}{y}$. Donc $x \mapsto
\frac{f(x)}x$ croît.

Sur-additivité pour $f(0) = 0$~: pour $x, y > 0$ (les cas avec une
variable nulle sont triviaux),
\[
f(x) = x\,\frac{f(x)}{x} \leq x\,\frac{f(x+y)}{x+y},
\qquad
f(y) \leq y\,\frac{f(x+y)}{x+y},
\]
par la monotonie qu'on vient de démontrer~; en additionnant on
obtient $f(x) + f(y) \leq f(x+y)$.
\end{solution}

\begin{solution}{exo:b2:realfun:10}
Pour $x, y \in I$ et $\lambda \in \intcc01$~: convexité de $f$, puis
monotonie de $g$, puis convexité de $g$~:
\[
g\bigl(f(\lambda x + (1{-}\lambda)y)\bigr)
\leq g\bigl(\lambda f(x) + (1{-}\lambda)f(y)\bigr)
\leq \lambda\,g(f(x)) + (1{-}\lambda)\,g(f(y)).
\]
Contre-exemple sans monotonie~: $g(t) = -t$ est convexe (affine)
mais décroissante, $f(x) = x^2$ est convexe, et $g \circ f = -x^2$
est strictement concave.
\end{solution}

\begin{solution}{exo:b2:realfun:11}
\emph{Inégalité de gauche~:} posons $m = \frac{a+b}2$ et prenons une
droite d'appui en $m$ (le~\cref{thm:b2:realfun:convexreg}~(2))~:
$f(t) \geq f(m) + \mu(t - m)$ pour tout $t \in \intcc ab$. En
intégrant sur $\intcc{a}{b}$~: le terme linéaire s'intègre en
$\mu\int_a^b(t - m)\dd t = 0$ (symétrie autour de $m$), donc
$\int_a^b f \geq (b - a)f(m)$.

\emph{Inégalité de droite~:} sur $\intcc ab$, la convexité majore $f$
par sa corde~: $f(t) \leq f(a) + \frac{f(b) - f(a)}{b - a}(t - a)$.
En intégrant~: $\int_a^b f \leq (b-a)f(a) + \frac{f(b) - f(a)}{b -
a}\cdot\frac{(b-a)^2}2 = (b - a)\,\frac{f(a) + f(b)}2$. Diviser par
$b - a$.
\end{solution}

\begin{solution}{exo:b2:realfun:12}
Soit $a - \delta < a \leq x < y \leq b < b + \delta$, tous dans $I$.
Deux applications de l'inégalité des pentes
(le~\cref{lem:b2:realfun:slopes}), d'abord à $a - \delta < a \leq x <
y$, puis à $x < y \leq b < b + \delta$~:
\[
\frac{f(a) - f(a - \delta)}{\delta}
\leq \frac{f(y) - f(x)}{y - x}
\leq \frac{f(b + \delta) - f(b)}{\delta}
\]
(les pentes des cordes croissent quand les deux extrémités se
déplacent vers la droite). Ainsi toute pente de corde dans
$\intcc ab$ est piégée entre deux nombres fixes, et
\[
\abs{f(y) - f(x)} \leq K\,\abs{y - x},
\qquad
K = \max\Bigl(\Bigl|\frac{f(a) - f(a-\delta)}{\delta}\Bigr|,
\Bigl|\frac{f(b+\delta) - f(b)}{\delta}\Bigr|\Bigr):
\]
$f$ est \omterm{def:b2:metric:continuity}{lipschitzienne} sur $\intcc ab$. Tout point de l'\omterm{def:b2:metric:topology}{ouvert} $I$
possède un tel segment-avec-marge autour de lui~: localement
\omterm{def:b2:metric:continuity}{lipschitzienne}, donc (de nouveau) \omterm{def:b2:metric:continuity}{continue} sur $I$.
\end{solution}

\begin{solution}{pb:b2:realfun:1}
\textbf{1.} \emph{Convexe $\Rightarrow$ $f'$ croissante~:} pour $a <
b$, l'inégalité des pentes donne, pour $h > 0$ petit, $\frac{f(a+h)
- f(a)}h \leq \frac{f(b) - f(a)}{b-a} \leq \frac{f(b) -
f(b-h)}{h}$~; en faisant $h \to 0$~: $f'(a) \leq \frac{f(b) -
f(a)}{b - a} \leq f'(b)$. \emph{Réciproquement}, si $f'$ croît et
$x < y < z$~: le théorème des accroissements finis donne $c_1 \in
\intoo{x}{y}$, $c_2 \in \intoo yz$ avec
\[
\frac{f(y) - f(x)}{y - x} = f'(c_1) \leq f'(c_2)
= \frac{f(z) - f(y)}{z - y},
\]
et cette inégalité des pentes à trois points, appliquée avec $y =
\lambda x + (1 - \lambda)z$, se réarrange en l'inégalité de
convexité. Pour le cas $C^2$~: $f'' \geq 0$ si et seulement si $f'$
croît.

\textbf{2.} Si $f'' > 0$, $f'$ est strictement croissante, et le
calcul par accroissements finis ci-dessus donne une inégalité
\emph{stricte} entre les deux pentes de corde~: stricte convexité.
\emph{Appui strict~:} en un point intérieur $a$ de pente d'appui $m$,
si $f(x_0) = f(a) + m(x_0 - a)$ pour un $x_0 \neq a$, alors sur le
segment de $a$ à $x_0$ la droite d'appui et la corde coïncident, et
la stricte convexité au milieu donne $f\bigl(\frac{a +
x_0}2\bigr) < f(a) + m\,\frac{x_0 - a}2$, ce qui contredit
l'inégalité d'appui. Donc $f(x) > f(a) + m(x - a)$ pour tout $x \neq
a$. \emph{Jensen strict~:} avec $a = \sum\lambda_ix_i$, en moyennant
les inégalités d'appui on obtient $\sum\lambda_if(x_i) \geq f(a)$,
avec égalité si et seulement si chaque terme est une égalité,
c'est-à-dire si et seulement si tout $x_i = a$.

\textbf{3.} $(-\ln)'' = \frac1{t^2} > 0$~; $(t^r)'' = r(r -
1)t^{r-2}$, positive pour $r > 1$, négative pour $0 < r < 1$~;
$\exp'' = \exp > 0$. Toutes strictes par la question 2.

\textbf{4.} Les cas $ab = 0$ sont triviaux. Pour $a, b > 0$,
concavité de $\ln$ aux points $a^p, b^q$ avec les poids
$\frac1p, \frac1q$~:
\[
\ln\Bigl(\frac{a^p}p + \frac{b^q}q\Bigr) \geq
\frac1p\ln(a^p) + \frac1q\ln(b^q) = \ln(ab),
\]
et $\ln$ croît~: $ab \leq \frac{a^p}p + \frac{b^q}q$.
Égalité si et seulement si les deux points coïncident (concavité
stricte)~: $a^p = b^q$.

\textbf{5.} Concavité de $\ln$ avec les poids $\lambda_i$~:
$\ln\bigl(\sum\lambda_ix_i\bigr) \geq \sum\lambda_i\ln x_i =
\ln\prod x_i^{\lambda_i}$~; exponentier. Égalité si et seulement si
tous les $x_i$ égaux (question 2). La voie de l'\cref{exo:b2:realfun:6}
obtenait la même inégalité comme limite de moyennes de puissances~;
ici c'est une seule application de Jensen --- la boîte à outils a de
la redondance intégrée.

\textbf{6.} Si $a = 0$ ou $b = 0$ l'inégalité est triviale.
Normalisons~: en remplaçant $a$ par $a/\norm a_p$ et $b$ par $b/\norm
b_q$, on peut supposer $\norm a_p = \norm b_q = 1$ et il faut montrer
$\sum\abs{a_ib_i} \leq 1$. Young terme à terme~:
\[
\sum_i\abs{a_i}\abs{b_i} \leq
\sum_i\Bigl(\frac{\abs{a_i}^p}{p} +
\frac{\abs{b_i}^q}{q}\Bigr) = \frac1p + \frac1q = 1 .
\]
Égalité si et seulement si chaque inégalité de Young est serrée~:
$\abs{a_i}^p = \abs{b_i}^q$ pour tout $i$ --- après avoir défait la
normalisation, $(\abs{a_i}^p)$ proportionnel à $(\abs{b_i}^q)$.

\textbf{7.} $p = q = 2$ est Cauchy--Schwarz avec le même cas
d'égalité (proportionnalité). Cas extrême~: $\sum\abs{a_ib_i}
\leq \bigl(\max_i\abs{b_i}\bigr)\sum_i\abs{a_i} =
\norm a_1\norm b_\infty$, immédiat terme à terme.

\textbf{8.} Pour $p = 1$ c'est l'inégalité triangulaire terme à
terme. Pour $p > 1$, avec $q$ conjugué~:
\[
\norm{a+b}_p^p = \sum_i\abs{a_i + b_i}^p
\leq \sum_i\abs{a_i+b_i}^{p-1}\abs{a_i} +
\sum_i\abs{a_i+b_i}^{p-1}\abs{b_i},
\]
et Hölder sur chaque somme, en notant $(p - 1)q = p$~:
\[
\sum_i\abs{a_i+b_i}^{p-1}\abs{a_i} \leq
\Bigl(\sum_i\abs{a_i+b_i}^{p}\Bigr)^{1/q}\norm a_p
= \norm{a + b}_p^{p/q}\,\norm a_p ,
\]
de même avec $b$. Donc $\norm{a+b}_p^p \leq \norm{a +
b}_p^{p/q}\bigl(\norm a_p + \norm b_p\bigr)$~; si $a + b \neq 0$,
diviser par $\norm{a+b}_p^{p/q}$ et utiliser $p - \frac pq = 1$. Avec
l'homogénéité et la séparation (claires), $\norm\cdot_p$ est une
\omterm{def:b2:nvs:norm}{norme} sur $\R^n$.

\textbf{9.} Pour $f, g$ \omterm{def:b2:metric:continuity}{continues} sur $\intcc ab$~: Hölder
\[
\int_a^b\abs{fg} \leq \Bigl(\int_a^b\abs
f^p\Bigr)^{1/p}\Bigl(\int_a^b\abs g^q\Bigr)^{1/q}
\]
par la même normalisation plus Young ponctuel, intégrée~; et
Minkowski $\norm{f + g}_p \leq \norm f_p + \norm g_p$ par le même
découpage, Hölder sur chaque morceau. La séparation de la \omterm{def:b2:nvs:norm}{norme}
utilise la stricte positivité~: un $\abs f^p$ \omterm{def:b2:metric:continuity}{continu} d'intégrale
nulle s'annule identiquement (volume de première année).

\textbf{10.} \emph{Monotonie~:} on peut supposer $\norm a_p = 1$~;
alors chaque $\abs{a_i} \leq 1$, donc $\abs{a_i}^q \leq \abs{a_i}^p$
et $\norm a_q^q \leq 1$~: $\norm a_q \leq 1 = \norm a_p$.
L'égalité exige $\abs{a_i}^q = \abs{a_i}^p$ pour tout $i$,
c'est-à-dire chaque $\abs{a_i} \in \{0, 1\}$~; avec
$\sum\abs{a_i}^p = 1$ cela ne laisse exactement qu'une coordonnée de
module $1$~: égalité si et seulement si $a$ a au plus une coordonnée
non nulle. \emph{Limite~:} $\norm a_\infty \leq
\norm a_p \leq n^{1/p}\norm a_\infty$, et $n^{1/p} \to 1$.
\emph{Comparaison inverse~:} Hölder avec les exposants $\frac
qp$ et son conjugué $\frac{q}{q-p}$, appliqué à
$\abs{a_i}^p\cdot 1$~:
\[
\norm a_p^p = \sum_i\abs{a_i}^p\cdot 1 \leq
\Bigl(\sum_i\abs{a_i}^{q}\Bigr)^{p/q}\,n^{1 - p/q}
= \norm a_q^{p}\; n^{1-p/q},
\]
d'où $\norm a_p \leq n^{\frac1p - \frac1q}\norm a_q$, avec
égalité si et seulement si tous les $\abs{a_i}$ sont égaux (le cas
d'égalité de Hölder contre le vecteur constant).

\textbf{11.} Écrivons $\abs{a_i}^r = \abs{a_i}^{\theta
r}\,\abs{a_i}^{(1-\theta)r}$ et appliquons Hölder avec les exposants
conjugués $\frac{p}{\theta r}$ et
$\frac{q}{(1-\theta)r}$ (conjugués précisément parce que
$\frac{\theta r}p + \frac{(1-\theta)r}q = 1$)~:
\[
\norm a_r^r = \sum_i \abs{a_i}^{\theta r}\abs{a_i}^{(1-\theta)r}
\leq \Bigl(\sum_i\abs{a_i}^{p}\Bigr)^{\theta r/p}
\Bigl(\sum_i\abs{a_i}^{q}\Bigr)^{(1-\theta)r/q}
= \norm a_p^{\theta r}\,\norm a_q^{(1-\theta)r} .
\]
Prendre les racines $r$-ièmes~: les $p$-normes sont log-convexes en
$\frac1p$.

\textbf{12.} \emph{Les deux négatifs~:} si $p < q < 0$ alors $0 < -q
< -p$, et $M_{-q}(y) \leq M_{-p}(y)$ pour les exposants positifs (cas
du cours, l'\cref{ex:b2:realfun:powermeans}) appliqué à $y =
(1/x_i)$~; en inversant l'identité $M_p(x) = M_{-p}(1/x)^{-1}$
l'inégalité se renverse en $M_p(x) \leq M_q(x)$.
\emph{Pont~:} pour $q > 0$, la concavité de $\ln$ donne $\ln M_q =
\frac1q\ln\bigl(\sum\lambda_ix_i^q\bigr) \geq
\frac1q\sum\lambda_i\ln x_i^q = \ln M_0$~; pour $p < 0$, la même
concavité donne $\ln\bigl(\sum\lambda_ix_i^p\bigr) \geq
p\sum\lambda_i\ln x_i$, et en divisant par $p < 0$ cela s'inverse~:
$\ln M_p \leq \ln M_0$. Donc $M_p \leq M_0 \leq M_q$ dès que $p < 0 <
q$~: avec les deux cas de même signe, $M$ croît sur tout
$\R^*$ (et à travers $0$).

\textbf{13.} Soit $x_{\max} = \max x_i$, atteint en $i^*$. Pour
$p > 0$~:
\[
\lambda_{i^*}^{1/p}\,x_{\max} \leq M_p \leq x_{\max},
\]
et $\lambda_{i^*}^{1/p} \to 1$~: $M_p \to x_{\max}$. Pour $p \to
-\infty$~: $M_p(x) = M_{-p}(1/x)^{-1} \to
\bigl(\max_i\frac1{x_i}\bigr)^{-1} = \min_ix_i$.

\textbf{14.} Avec $\lambda_i = \frac1n$, la chaîne $M_{-\infty}
\leq M_{-1} \leq M_0 \leq M_1 \leq M_2 \leq M_{+\infty}$ s'écrit
\[
\min \leq \frac{n}{\sum\frac1{a_i}} \leq \Bigl(\prod
a_i\Bigr)^{1/n} \leq \frac{\sum a_i}{n} \leq
\sqrt{\frac{\sum a_i^2}{n}} \leq \max .
\]
La comparaison arithmético-harmonique ($M_{-1} \leq M_1$) se
réarrange directement en
$\bigl(\sum a_i\bigr)\bigl(\sum\frac1{a_i}\bigr) \geq n^2$.

\textbf{15.} Avec des poids égaux, $M_p(x) =
\bigl(\frac1n\sum\abs{x_i}^p\bigr)^{1/p} = n^{-1/p}\norm x_p$.
Quand $p$ croît, $\norm x_p$ décroît (question 10) mais le
normalisateur $n^{-1/p}$ croît plus vite, et le produit
croît (question 12)~: les moyennes font la moyenne, les \omterm{def:b2:nvs:norm}{normes}
accumulent, et le facteur $n^{-1/p}$ est exactement le taux de change
entre les deux conventions comptables.

\textbf{16.} Chaque maillon est une instance de Jensen strict
(question 2) avec les fonctions strictement convexes/concaves de la
question 3 ($t^{q/p}$, $\ln$), donc l'égalité à un maillon quelconque
force tous les $x_i$ égaux~; et $\min = M_p$ ou $M_p = \max$ force de
même toutes les valeurs égales à l'extremum commun. La chaîne est
stricte dès que deux $x_i$ diffèrent.

\textbf{17.} Appliquer Young (question 4) au couple
$\varepsilon^{1/p}a$ et $\varepsilon^{-1/p}b$~:
\[
ab = (\varepsilon^{1/p}a)(\varepsilon^{-1/p}b)
\leq \varepsilon\,\frac{a^p}p +
\varepsilon^{-q/p}\,\frac{b^q}q .
\]
Pour $p = q = 2$, en remplaçant $\varepsilon$ par $2\varepsilon$~:
$ab \leq \varepsilon a^2 + \frac{b^2}{4\varepsilon}$ --- l'inégalité
d'absorption~: un produit est échangé contre un petit multiple d'un
carré plus un grand multiple de l'autre.

\textbf{18.} Télescopage~:
\[
\prod_{k=1}^{n}c_k = \frac{\prod_{k=1}^n(k+1)^k}
{\prod_{k=1}^{n}k^{k-1}}
= \frac{2^1\,3^2\cdots(n+1)^n}{1^0\,2^1\cdots n^{n-1}}
= (n+1)^n,
\]
chaque facteur $(k+1)^k$ du numérateur s'annulant contre le terme
suivant du dénominateur. Inégalité arithmético-géométrique sur les
$n$ nombres $c_ka_k$~:
\[
(a_1\cdots a_n)^{1/n} =
\frac{\bigl(\prod_k c_ka_k\bigr)^{1/n}}{(n+1)}
\leq \frac{1}{n+1}\cdot\frac1n\sum_{k=1}^{n}c_ka_k .
\]

\textbf{19.} En sommant sur $n$ et en échangeant les deux sommations
(tous les termes positifs~: le~\cref{thm:b2:series:fubini})~:
\[
\sum_{n\geq1}(a_1\cdots a_n)^{1/n}
\leq \sum_{n\geq1}\frac{1}{n(n+1)}\sum_{k=1}^{n}c_ka_k
= \sum_{k\geq1}c_ka_k\sum_{n\geq k}\frac1{n(n+1)}
= \sum_{k\geq1}\frac{c_ka_k}{k},
\]
en utilisant le télescopage $\sum_{n\geq k}\bigl(\frac1n -
\frac1{n+1}\bigr) = \frac1k$. Enfin $\frac{c_k}k =
\frac{(k+1)^k}{k^k} = \bigl(1 + \frac1k\bigr)^k < \eu$
(suite croissante de limite $\eu$, volume de première année)~:
\[
\sum_{n\geq1}(a_1\cdots a_n)^{1/n} \leq
\eu\sum_{k\geq1}a_k :
\]
l'inégalité de Carleman. (La constante $\eu$ est optimale, bien
qu'on ne le démontre pas.)

\textbf{20.} Cauchy--Schwarz (question 9, $p = q = 2$) appliqué
à $\sqrt f$ et $\frac1{\sqrt f}$~:
\[
1 = \Bigl(\int_0^1\sqrt f\cdot\frac{1}{\sqrt f}\Bigr)^{2}
\leq \Bigl(\int_0^1 f\Bigr)\Bigl(\int_0^1\frac1f\Bigr).
\]
Égalité si et seulement si $\sqrt f$ et $\frac1{\sqrt f}$ sont
proportionnels, c'est-à-dire $f^2$ constant, c'est-à-dire $f$
constant ($f > 0$ \omterm{def:b2:metric:continuity}{continue}).

\textbf{21.} Soit $1 < p < \infty$, $\norm a_p = \norm b_p = 1$,
$a \neq b$, et supposons $\bigl\Vert\frac{a+b}2\bigr\Vert_p = 1$,
c'est-à-dire Minkowski est une égalité pour $a, b$. En retraçant la
démonstration de la question 8, l'égalité force l'égalité dans les
deux applications de Hölder et dans les inégalités triangulaires
terme à terme~: $(\abs{a_i}^p)$ et $(\abs{b_i}^p)$ tous deux
proportionnels à $(\abs{a_i + b_i}^p)$, et $a_i, b_i$ de même signe
--- d'où $b = ta$ pour un $t \geq 0$, et $\norm b_p = \norm a_p$
donne $t = 1$~: $b = a$, contradiction. Donc la $p$-sphère ne
contient aucun milieu de points distincts de la sphère~: aucun
segment. Pour $p = \infty$ dans $\R^2$~: tous les $(1, t)$,
$\abs t \leq 1$, sont sur la sphère unité --- un bord plat~; pour
$p = 1$~: le segment $(t, 1 - t)$, $t \in \intcc01$, y est.

\textbf{22.} Pour $a = 0$ les deux membres s'annulent. Sinon Hölder
majore tout $\sum a_ib_i$ par $\norm a_p\norm b_q \leq \norm
a_p$. Atteinte~: prenons
\[
b_i = \frac{\operatorname{sign}(a_i)\,\abs{a_i}^{p-1}}
{\norm a_p^{p/q}} :
\qquad
\norm b_q^q = \frac{\sum_i\abs{a_i}^{(p-1)q}}{\norm a_p^{p}}
= \frac{\norm a_p^p}{\norm a_p^p} = 1,
\quad
\sum_ia_ib_i = \frac{\norm a_p^p}{\norm a_p^{p/q}} =
\norm a_p ,
\]
en utilisant $(p-1)q = p$ et $p - \frac pq = 1$. Donc la borne
supérieure est un maximum, égal à $\norm a_p$~: chaque $p$-norme est
la \omterm{def:b2:nvs:norm}{norme} duale de sa conjuguée --- le germe de la dualité
$L^p$--$L^q$.

\textbf{23.} $\E[X^r] = \sum_i\lambda_ix_i^r$, donc
$\E[X^r]^{1/r} = M_r(x; \lambda)$, croissante en $r$ par la
question 12 (et à travers $r \to 0, \pm\infty$ par les questions
12--13)~: l'inégalité des moments de Lyapunov, purement un énoncé sur
les moyennes de puissances pondérées. Elle revient pour de véritables
variables aléatoires au~\cref{ch:b2:randomvar}.

\textbf{24.} (i) Moyennes de puissances $M_1 \leq M_3$ avec des
poids égaux~: $\frac{a+b+c}3 \leq \bigl(\frac{a^3+b^3+c^3}3\bigr)^{1/3}$~;
élever au cube et multiplier par $3$~: $a^3 + b^3 + c^3 \geq
\frac{(a+b+c)^3}9$. (ii) Cauchy--Schwarz contre le vecteur
constant~: $\sum_i\sqrt{x_i}\cdot1 \leq
\bigl(\sum_ix_i\bigr)^{1/2}n^{1/2}$~; élever au carré.

\textbf{25.} La définition par corde livre le lemme des pentes par
un seul réarrangement algébrique~; les pentes resserrées en un point
produisent des dérivées latérales et des droites d'appui, dont la
moyenne pondérée est Jensen. Appliqué à $-\ln$, Jensen devient Young,
qui sommé contre des vecteurs normalisés est Hölder, qui découpé et
réabsorbé est Minkowski --- et les $p$-normes du~\cref{ch:b2:nvs}
naissent, avec leur dualité (question 22) et leur géométrie
(question 21). Jensen appliqué le long de l'échelle des puissances
enchaîne toutes les moyennes de $\min$ à $\max$ (questions 12--14),
qui lues sur des variables aléatoires est l'inégalité des moments
(question 23). Et l'inégalité arithmético-géométrique, pondérée par
une astuce de télescopage, livre la borne de Carleman avec sa
constante irréductible $\eu$ (questions 18--19). Sommets~:
Hölder--Minkowski, et Carleman. Destination~: les espaces $L^p$
du volume de troisième année, dont les axiomes fondateurs sont
exactement les questions 6 et 8 avec des intégrales à la place des
sommes.
\end{solution}
